\section{Introduction}\label{sec:introduction}

\subsection{Goal of the project}
The goal of this project is to set up an agent-based transport demand pipeline for one French department outside
the \^Ile-de-France region. The pipeline starts with a territorial analysis of the study area, generates and attracts
trips per zone, distributes them with a double-constrained gravity model, disaggregates the resulting flows into
individual trips with coordinates and departure times, and finally routes these trips on a road network extracted from
OpenStreetMap. The simulation does not aim to be correct; it shows that the tools can be chained and that analyses can
be performed on their outputs.
\todo{Rewrite the goal in the authors' own words, 4--6 sentences, and add a short overview figure of the pipeline if useful.}

\subsection{Study area}
\todo{Name the chosen department (code and name), its location, its main urban centres and the reasons for the choice
(size, availability of data, a mix of urban and rural municipalities, \dots). The detailed census figures belong to
\cref{sec:study-area}.}

\subsection{Data sources}\label{sec:data-sources}
The following data sets are used (links, file names and known pitfalls are listed in \code{data/DATA\_SOURCES.md} of the
repository):
\begin{itemize}
  \item the INSEE detailed census, ``Individus localis\'es au canton-ou-ville'' \cite{insee-census-detailed};
  \item the INSEE aggregated census, ``Population en 2019 -- IRIS'' \cite{insee-census-aggregated};
  \item the IGN CONTOURS-IRIS boundaries, 2021 edition \cite{ign-contours-iris};
  \item the URSSAF number of employees per municipality and economic activity (APE) \cite{urssaf-employment};
  \item the OpenStreetMap extract of the region from Geofabrik \cite{geofabrik-osm}.
\end{itemize}
\todo{Add vintages, access dates and the exact files used. Mention that the detailed census is localised at the level of
``canton-ou-ville'' and not at municipality level (see \cref{sec:study-area}).}

\subsection{Tools}
The analysis is written in Python (geopandas \cite{geopandas}, OSMnx \cite{osmnx}, pyrosm \cite{pyrosm}) in Jupyter
notebooks, supported by a small shared package \code{lsm} in \code{src/}. The report is written in \LaTeX{}.
\todo{State the exact versions (taken from \code{environment.yml}) and refer to \cref{sec:appendix-reproducibility}.}

\subsection{Team split}
\Cref{tab:team-split} summarises who is responsible for which exercise. The only interface between the two parts is the
file \code{data/processed/zones.parquet} (zone identifier, name, geometry, \Oprime{}, \Dj{}), described in
\cref{sec:appendix-reproducibility}.

\begin{table}[H]
  \centering
  \caption{Team split. Points are those given in the assignment.}
  \label{tab:team-split}
  \begin{tabularx}{\textwidth}{@{}llXr@{}}
    \toprule
    Person & Name & Exercises & Points \\
    \midrule
    A & \todo{Name One} & 1.1 Study area, 1.2 and 1.3 Territorial analysis, 2.1 Trip production, 2.2 Trip attraction & 12 \\
    B & \todo{Name Two} & 2.3 Trip distribution, 3.1 Disaggregation, 3.2 Routing & 10 \\
    \bottomrule
  \end{tabularx}
\end{table}

\subsection{Structure of this report}
Each of \crefrange{sec:study-area}{sec:routing} answers one exercise of the assignment, one subsection per question, and
ends with the problems that were encountered. \Cref{sec:discussion} discusses limitations and consistency checks across
the pipeline.
